\documentclass[10pt,a4paper]{amsart}
\usepackage[margin=19mm]{geometry}
\usepackage{amsmath,amssymb,mathtools}
\usepackage[scr=rsfs,cal=euler]{mathalfa}
\usepackage{xfrac}
\usepackage[unicode,hidelinks,bookmarks=false]{hyperref}
\newcommand{\ie}{i.e.\ }
\newcommand{\cf}{cf.\ }
\newcommand{\resp}{respectively\ }
\newcommand{\Subsection}{\subsection}
\providecommand{\nobreakdash}{\nobreak\hbox{-}}
\setlength{\emergencystretch}{2em}
\multlinegap0pt
\usepackage{mathtools}
\usepackage{amssymb}
\usepackage{tikz}
\newtheorem{lemma}{Lemma}
\title{Klein bottle and optimal systolic inequality for nonpositively curved surfaces}
\author{Mikhail G. Katz, Stephane Sabourau}
\date{Source: arXiv:2609.06705v1 (CC BY 4.0)}
\begin{document}
\maketitle

\noindent\emph{Source and licence.} Klein bottle and optimal systolic inequality for nonpositively curved surfaces. Version arXiv:2609.06705v1 (CC BY 4.0), distributed by arXiv under the Creative Commons Attribution 4.0 International licence, http://creativecommons.org/licenses/by/4.0/, as recorded in the arXiv licence metadata for this exact version and shown on the abstract page https://arxiv.org/abs/2609.06705v1. Full licence notice: \texttt{licenses/arxiv-2609.06705-CC-BY-4.0.txt}.\par\smallskip

\noindent\textit{Editorial scope.} Counted unit: the article's lemma lemma (source0.tex lines 633--635). The article's complete original proof of it is delivered with it (source0.tex lines 637--660, one \texttt{proof} environment of 24 lines). No mathematics is written, repaired or re-derived: the statement and the proof are the article's own lines, in the article's own notation, and no retained line is substituted or re-flowed.  Retained uncounted as the source-local context the proof needs: lines 564--566.  Nothing was omitted from any retained span, so the package carries no omission record.  No retained line contains a citation, so the wrapper carries no bibliography and no bibliography entry.  No retained line contains a figure environment, an image-inclusion command or drawing code, so no image is delivered and the package declares no asset dependency.  Editorial formatting only, in addition: an AMS \texttt{amsart} preamble that restates the article's own declarations and macros used by the retained lines, namely the environments \texttt{lemma} on separate counters, together with the standard packages those lines need.  The retained environments print in this document as Lemma 1 in order of appearance, whereas the article prints the same items with the numbering of its own full document.  The complete native source of the article as distributed in the arXiv e-print for this exact version is bundled unchanged as \texttt{sources/209562/source0.tex}.
\begin{lemma} \label{lem:noncontracting}
The map $\Phi_\theta: \mathcal{T}_\theta \setminus D_\theta \to \bar{M} \setminus D_0$ is locally distance-nondecreasing.
\end{lemma}

\begin{lemma}
The image of~$Q$ by the exponential-type map lies within the Voronoi cell~$V_0$.
\end{lemma}

\begin{proof}
Consider a point~$\omega \in Q$ and its image $y=\Phi_\theta(\omega)$ in~$\bar{M}$.
By construction, the point~$\omega$ lies in the region containing the origin~$\omega_0$ and bounded by the equidistant lines~$\Delta_i^\theta$ in~$\mathcal{T}_\theta$.
Therefore,
\[
d_{\bar{M}}(x_0,y) = d_{\mathcal{T}_\theta}(\omega_0,\omega) \leq d_{\mathcal{T}_\theta}(\omega_i,\omega).
\]

Suppose the minimizing arc~$\alpha \subseteq \bar{M}$ between~$x_i$ and~$y$ avoids~$D_0$.
Then its preimage~$\Phi^{-1}_\theta(\alpha) \subseteq \mathcal{T}_\theta$ is an arc between~$\omega_i$ and~$\omega$, which is no longer than~$\alpha$ by Lemma~\ref{lem:noncontracting}.
Thus,
\[
d_{\bar{M}}(x_0,y) \leq d_{\mathcal{T}_\theta}(\omega_i,\omega) \leq d_{\bar{M}}(x_i,y).
\]
It follows that $y$ lies in the Voronoi cell~$V_0$ centered around~$x_0$.

Suppose the minimizing arc~$\alpha$ intersects~$D_0$.
Since $D_0$ is contained in the Voronoi cell~$V_0$, a point~$z$ of this arc is closer to~$x_0$ than~$x_i$.
Hence,
\[
d_{\bar{M}}(x_0,y) \leq d_{\bar{M}}(x_0,z) + d_{\bar{M}}(z,y) \leq d_{\bar{M}}(x_i,z) + d_{\bar{M}}(z,y) = d_{\bar{M}}(x_i,y).
\]
It follows as before that $y$ lies in~$V_0$.
\end{proof}

\end{document}
